% Bagean: metodhe
% Bab: induksi
% Pérangan: pambuka

\documentclass[../../../include/open-logic-section]{subfiles}

\begin{document}

\olfileid{mth}{ind}{int}

\olsection{Pambuka}

Induksi iku tèknik bukti kang wigati lan digunakaké,
kanthi wujud kang béda-béda, ing meh kabèh pérangan
logika, èlmu komputer téorètis, lan matématika.
Tèknik iki dibutuhaké kanggo mbuktekaké akèh asil
ing logika.

Induksi kerep dibédakaké saka deduksi lan diterangaké
minangka inferensi saka bab khusus menyang bab umum.
Upamané, yèn kita mirsani akèh watu jamrud ijo lan
ora nemoni jamrud kang ora ijo, kita bisa nyimpulaké
yèn kabèh jamrud iku ijo. Iki inferensi induktif,
amarga mangkat saka akèh kasus tartamtu, kaya
jamrud iki ijo lan jamrud kuwi ijo, menyang pratelan
umum yèn kabèh jamrud ijo. Induksi
\emph{matématis} uga ngasilaké pratelan umum,
nanging jinis inferensiné béda banget saka
``induksi prasaja'' mau.

Sacara kasar, bukti induktif ing matématika
nyimpulaké yèn kabèh obyèk matématis saka
sawijining jinis nduwèni sipat tartamtu.
Ing kasus kang paling prasaja, obyèk mau
yaiku bilangan asli. Bukti induktif banjur
netepaké yèn kabèh bilangan asli nduwèni
sipat tartamtu kanthi nuduhaké yèn
\begin{enumerate}
    \item $0$ nduwèni sipat mau, lan
    \item saben bilangan~$k$ kang nduwèni sipat mau,
    $k+1$ uga nduwèni.
\end{enumerate}
Induksi ing bilangan asli asring uga bisa
mbuktekaké pratelan umum babagan obyèk
matématis kang bisa diwènèhi cacah.
Upamané, saben himpunan winates nduwèni
cacah unsur winates~$n$. Yèn induksi
bisa nuduhaké yèn saben bilangan~$n$
nduwèni sipat ``kabèh himpunan winates
mawa ukuran~$n$ iku \dots'', kita wis
mbuktekaké sawijining bab babagan kabèh
himpunan winates.

Induksi uga bisa digeneralisasi marang obyèk
matématis kang \emph{didhéfinisikaké kanthi induktif}.
Upamané, ekspresi basa formal, kaya ekspresi
ing logika orde kapisan, didhéfinisikaké
kanthi induktif. \emph{Induksi struktural}
iku cara kanggo mbuktekaké asil babagan
kabèh ekspresi kaya mangkono. Mliginé,
induksi struktural migunani banget lan
digunakaké ing endi-endi ing logika.

\end{document}
